\subsection{带基点同伦}

设 X 与 Y 是拓扑空间，x 是 X 的一点。若同伦 \(\sigma\colon X \times I \to Y\) 在 \(\{x\}\) 上是固定的，即若从 I 到 Y 的映射 \(t \mapsto \sigma(x, t)\) 是常值的，则称它在 x 处带基点。两个在 x 处带基点的同伦的拼接在 x 处带基点。在 x 处带基点的同伦的起点与终点在 x 处取同一值 y；因此它们是从 \((X, x)\) 到 \((Y, y)\) 的带基点连续映射。

设 \((X,x)\) 与 \((Y,y)\) 是带基点拓扑空间（I, p. 120, 定义 1）。关系“f 由在 x 处带基点的同伦连接到 g”是集合 \(\mathcal{C}((X,x);(Y,y))\) 上的一个等价关系，称为带基点同伦关系。\(\mathcal{C}((X,x);(Y,y))\) 对该等价关系的商集记作 \([(X,x);(Y,y)]\)。它的元素称为从 \((X,x)\) 到 \((Y,y)\) 的带基点连续映射的带基点同伦类。

\pandocbounded{\includegraphics[keepaspectratio]{images/part002_b3045f29.jpg}}

集合 \(\mathcal{C}((\mathrm{X},x);(\mathrm{Y},y))\) 是 \(\mathcal{C}(\mathrm{X};\mathrm{Y})\) 的一个子集。注意，带基点同伦是比同伦在 \(\mathcal{C}(\mathrm{X};\mathrm{Y})\) 上诱导的关系更细的等价关系。一般而言它是一个严格更细的关系（见 III, p. 321, 习题 1 与 III, p. 234, 例 3）。

再设 \((Z,z)\) 是带基点拓扑空间。设 f 与 \(f'\) 是从 \((X,x)\) 到 \((Y,y)\) 的带基点连续映射，g 与 \(g'\) 是从 \((Y,y)\) 到 \((Z,z)\) 的带基点连续映射。若 f 与 \(f'\) 由在 x 处带基点的同伦 \(\sigma\) 相联系，且 g 与 \(g'\) 由在 y 处带基点的同伦 \(\tau\) 相联系，则 \(g \circ f\) 与 \(g' \circ f'\) 由同伦 \(\theta: X \times I \to Z\)，\((u,t) \mapsto \tau(\sigma(u,t),t)\) 相连接，该同伦在 x 处带基点。与上文一样，由此可以定义从 \([(X,x);(Y,y)] \times [(Y,y);(Z,z)]\) 到 \([(X,x);(Z,z)]\) 的合成映射。我们让读者自己陈述并证明关于带基点映射与带基点同伦的、与命题 3 及其推论相仿的命题。

